\begin{pastexam}[label={px:df3-105-ds-02}]{地方三等資訊處理 105 年 第二題}
（申論；15 分）

分別給定矩陣 A、B、C 與 D 的大小為 2×4、4×3、3×5 和 5×1：（每小題 5 分，共 15 分）
\begin{enumerate}[leftmargin=2.6em, labelsep=0.4em]
\item[㈠] 共有幾種加括號的方法？
\item[㈡] 例如(AB)(CD)，共需多少次乘法？
\item[㈢] 求出三者乘積之最有效的方式為何？
\end{enumerate}
\end{pastexam}

\begin{pastsoln}
㈢的「三者」應是筆誤，題目給的是四個矩陣，以下求 $ABCD$ 四者乘積最省的方式。 和
歷屆試題~\ref{px:gk3-110-ds-01} 是同一類題目（矩陣鏈乘，資結書第~18.2~節），這裡鏈比較短。

大小序列 $\langle p_0, \ldots, p_4 \rangle = \langle 2, 4, 3, 5, 1 \rangle$；$p \times q$ 乘 $q \times r$ 要
$pqr$ 次純量乘法。

\textbf{㈠} $4$ 個矩陣加滿括號的方法有卡特蘭數 $C_3 = 5$ 種：
\[
A(B(CD)),\quad A((BC)D),\quad (AB)(CD),\quad (A(BC))D,\quad ((AB)C)D .
\]
（依最後一次乘法的切點分類：$C_0 C_2 + C_1 C_1 + C_2 C_0 = 2 + 1 + 2 = 5$。）

\textbf{㈡} $(AB)(CD)$：
\begin{itemize}
\item $AB$：$(2 \times 4)(4 \times 3)$，$2 \cdot 4 \cdot 3 = 24$ 次，得 $2 \times 3$；
\item $CD$：$(3 \times 5)(5 \times 1)$，$3 \cdot 5 \cdot 1 = 15$ 次，得 $3 \times 1$；
\item $(AB)(CD)$：$(2 \times 3)(3 \times 1)$，$2 \cdot 3 \cdot 1 = 6$ 次。
\end{itemize}
共 $24 + 15 + 6 = 45$ 次乘法。

\textbf{㈢} 用動態規劃，$m[i,j] = \min_{i \le k < j}\{m[i,k] + m[k+1,j] + p_{i-1} p_k p_j\}$，$m[i,i] = 0$
（$A_1, \ldots, A_4$ 依序是 A、B、C、D）：
\begin{itemize}
\item 長度 $2$：$m[1,2] = 2 \cdot 4 \cdot 3 = 24$，$m[2,3] = 4 \cdot 3 \cdot 5 = 60$，$m[3,4] = 3 \cdot 5 \cdot 1 = 15$。
\item $m[1,3] = \min\{0 + 60 + 2 \cdot 4 \cdot 5,\; 24 + 0 + 2 \cdot 3 \cdot 5\} = \min\{100, 54\} = 54$（$k = 2$）。
\item $m[2,4] = \min\{0 + 15 + 4 \cdot 3 \cdot 1,\; 60 + 0 + 4 \cdot 5 \cdot 1\} = \min\{27, 80\} = 27$（$k = 2$）。
\item $m[1,4] = \min\{0 + 27 + 2 \cdot 4 \cdot 1,\; 24 + 15 + 2 \cdot 3 \cdot 1,\; 54 + 0 + 2 \cdot 5 \cdot 1\}
= \min\{35, 45, 64\} = 35$（$k = 1$）。
\end{itemize}
$m[1,4]$ 在 $k = 1$ 取到，切成 $A \mid BCD$；$m[2,4]$ 在 $k = 2$，切成 $B \mid CD$。 最有效的方式是
\[
A(B(CD)),
\]
共 $35$ 次乘法：$CD$ 要 $15$ 次得 $3 \times 1$，$B(CD)$ 要 $4 \cdot 3 \cdot 1 = 12$ 次得 $4 \times 1$，
$A(B(CD))$ 要 $2 \cdot 4 \cdot 1 = 8$ 次，$15 + 12 + 8 = 35$。 五種順序全部算出來是：$A(B(CD))$ $35$、
$(AB)(CD)$ $45$、$((AB)C)D$ $64$、$A((BC)D)$ $88$、$(A(BC))D$ $110$，最少的確是 $35$。 直觀的理由：
D 只有 $1$ 行，從右邊先乘，每一步的結果都只有 $1$ 行，乘法次數都很小。
\end{pastsoln}
